\section{Reduction to spherical moments}\label{sec:sphere}

Contraction kills the image of each \(\xi_i-\partial_{z_i}\). Applying
this operator to \(\xi^{\alpha-e_i}z^\beta\) replaces one outer factor
\(\xi_i\) by \(\partial_{z_i}\) modulo its image. Induction on the
\(\xi\)-degree shows that every \(H-\mathcal E(H)\) belongs to the sum
of the two images, proving the kernel identity.

Introduce the four balanced bilinears
\[
 R=\xi_1z_1+\xi_2z_2,\quad Z=\xi_1z_2,\quad
 W=2\xi_2z_1,\quad T=\xi_1z_1-\xi_2z_2,
 \qquad T^2+2ZW=R^2.
\]
They generate the balanced subring, with the displayed relation.
Generation follows by pairing one \(\xi_i\) with one \(z_j\). The normal
monomials modulo the relation in homogeneous degree \(d\) number
\((d+1)^2\), the dimension of the original bidegree-\((d,d)\) space,
so there is no further relation. Restrict at \(R=1\), writing
\((x,y,t)=(Z,W,T)\) there. The sphere algebra and its trace are
\[
 A_K=K[x,y,t]/(t^2+2xy-1),\qquad
 \Lambda(h)=\operatorname{CT}_x\frac12\int_{-1}^1
    h\biggl(x,\frac{1-t^2}{2x},t\biggr)\,dt.
\]
The integral is the rational polynomial rule with value \(1/(j+1)\)
on \(t^j\) for even \(j\), and zero for odd \(j\), including the factor
\(1/2\). If \(H\) has bidegree \((d,d)\) and sphere restriction \(h\), then
\begin{equation}\label{eq:spectrum-contraction}
 \mathcal E(H)=(d+1)!\,\Lambda(h).
\end{equation}
To check this, take an original monomial with \(\xi_1\)-exponent \(i\)
and \(z_1\)-exponent \(j\). Both sides vanish for \(i\ne j\) by angular
weight. For \(i=j\) the sphere integral is the beta integral
\(i!(d-i)!/(d+1)!\), and its contraction is \(i!(d-i)!\).

The least sphere degree of \(h\) is \(d\) exactly when \(R\) does not
divide \(H\). Indeed homogenizing the normal monomials \(x^at^b\),
\(y^ct^b\) with \(c>0\) recovers \(H\), and each has filtration degree
\(a+b\) or \(c+b\), respectively. We use \emph{pure} for the condition
\(\Lambda(h^m)=0\) for every \(m\geq 1\) and \emph{safe} for the assertion
that \(\Lambda(gh^m)=0\) eventually for each fixed \(g\in A_K\).

Sphere safety gives the required vanishing against arbitrary original
polynomial markers. For a bihomogeneous marker \(Q\) of bidegree
\((a,b)\), contraction against \(H^m\) is zero if \(a>b\). For \(b\ge a\),
its output has degree \(b-a\) in \(z\), and for every \(|\gamma|=b-a\),
\[
 \partial_z^\gamma\mathcal E(QH^m)
   =\mathcal E(\xi^\gamma QH^m).
\]
The marker on the right is balanced. There are finitely many such
derivatives, so their eventual vanishing by
\eqref{eq:spectrum-contraction} forces eventual vanishing of the
output. A general marker is the sum of finitely many bihomogeneous
ones. Conversely any fixed sphere marker with infinitely many nonzero
mixed moments lifts to a balanced marker by homogenization.

\section{Quadratic height and the degrees below four}\label{sec:quadratic-height}

We first treat phases whose Laurent expression after eliminating \(y\)
has degree at most two in \(t\). A coefficient of \(x^{-j}\), \(j>0\),
is divisible by \((1-t^2)^j\), so only a simple pole at \(x=0\) is possible.
With no pole, the angular core must vanish for purity by the interval
polynomial theorem \cite[Theorem~2.1]{muger-tuset}. All nonzero angular exponents are then positive,
giving safety. With a pole, scale the phase and write
\[
 h=x^{-1}(A-Bt^2+Ct),\qquad A,B,C\in K[x],\quad A(0)=B(0)=1,\quad C(0)=0.
\]
Put \(D=A+C^2/(4B)\) and \(M=\sqrt{D/B}\), with \(M(0)=1\).
In completing the square, set \(t=C/(2B)+Mz\). The endpoints in \(z\)
are congruent to \(\pm1\) modulo \(x\). The integrand \((1-z^2)^m\)
has contact order \(m\) there, so returning the endpoints to \(\pm1\)
changes its integral by a multiple of \(x^{m+1}\). Therefore
\[
 \Lambda(h^m)=\frac12\int_{-1}^1(1-z^2)^m\,dz\;[x^m]D^mM.
\]
The factor given by the integral is a nonzero rational number.
The conditions \([x^m]D^mM=0\) determine each successive coefficient of
\(M\) since \(D(0)=1\). Differentiating \(D^m\) shows that the unique
solution is \(M=1-xD'/D\). In particular
\begin{equation}\label{eq:quintic-quadratic-height}
 B=\frac{D^3}{(D-xD')^2}.
\end{equation}

Suppose the phase has least sphere degree at most five, so \(\deg B\le4\).
At any finite zero or pole \(\alpha\) of \(D\) of order \(r\ne0\),
we have \(\alpha\ne0\), \(\operatorname{ord}_\alpha(D-xD')=r-1\), and
\(\operatorname{ord}_\alpha B=r+2\). There are no other finite zeros of
\(B\) by \eqref{eq:quintic-quadratic-height}. If \(n\) is the degree of
\(D\) at infinity and \(s\) counts its distinct finite zeros and poles
over an algebraic closure, summing valuations gives \(\deg B=n+2s\).
When \(n\ne1\), the leading coefficient of \(D-xD'\) is \(1-n\)
times that of \(D\), hence \(\deg B=n\) and \(D\) is constant.
When \(n=1\), the bound \(\deg B\le4\) forces \(s=1\), so \(D\) is
linear and \(B=D^3\). Then \(B\) divides \(C^2\), and \(C=xD^2P\)
for a polynomial \(P\). Since \(A=D-x^2DP^2/4\) has degree at most six,
\(P\) has degree at most one. The phase has exact degree four, because
its \(x^{-1}Bt^2\) term has degree four and its other terms have degree
at most four.

In the constant case \(D=B=1\). With \(c=C/x\), the phase is
\(2y+ct-xc^2/4\), the image of \(2y\) under the polynomial shear
\((x,t,y)\mapsto(x,t+xa,y-at-xa^2/2)\) with \(a=-c/2\).
The derivation \(xa\partial_t\) in the literal Laurent chart kills the
trace by integration at \(t=\pm1\). Each boundary value is a polynomial
in \(x\) since \(y=0\) there, and the extra factor \(x\) kills its constant
term. The shear is its finite exponential on each regular polynomial,
so both directions preserve the trace and safety. Thus in quadratic
height every pure phase of least degree at most five is safe or has
degree four.

Here is the exclusion in degrees \(d=1,2,3\), using the leading-form
test proved in Section~\ref{sec:leading-threshold}. In a pure phase the
binary leading form has a root of multiplicity greater than \(d\).
Move the root to the point \(v=0\) in the parametrization
\((x,y,t)=(u^2,-v^2/2,uv)\). Every degree-\(d\) normal monomial then
has negative angle \(i\). If the phase has no positive angle it is
safe by the interval theorem and one-sided support. Otherwise, use the
supporting-line and two-term tests of Section~\ref{sec:face-support} with
\(d\) in place of five. The allowed nodes have \(-d\le i\le d-1\) and
\(\max(0,-2i)\le j\), with \(i+j\le d\) for \(i<0\) and
\(i+j\le d-1\) otherwise. A crossing line \(j+\rho i=H\) has
\(\rho>1\) and \(H<d\).

Joining opposite-sign nodes in that range gives no line with three
or more nodes for \(d=1,2\). For \(d=3\) the only such lines are
\((\rho,H)=(2,2)\), with angles \(-1,0,1\), and
\((\rho,H)=(3/2,3/2)\), with angles \(-3,-1,1\). The former has
quadratic height after interchanging \(x,y\), so is safe by the
preceding argument. The latter has face \(Bx+Ayt+Cy^3\) with all
three coefficients nonzero. Its order-three cover has branch
exponents \(2,2,2\) when the face quadratic has distinct roots, so its
elliptic Jacobian supplies the ordinary-prime obstruction of
Section~\ref{sec:cyclic-face}. If the quadratic repeats, then
\(A^2+2BC=0\). The face is \(B\) times the leading part of
\(x\mapsto x-cyt-c^2y^3/2\) with \(c=-A/B\), in the shear fixing
\(y\) and sending \(t\) to \(t+cy^2\). Its inverse preserves the weights
\((3,1,2)\) on \((x,y,t)\) and removes the cubic part, leaving
degree at most two. This last phase is safe by the preceding
low-degree cases. Degree zero is immediate since a nonzero constant
has nonzero first moment. The sharper mixed-moment bound for cubics
is proved in \cite{cubicsphere}.

\section{Explicit primitive witnesses}\label{sec:witnesses}

The following construction uses the quartic endpoint mechanism of
Long \cite[Theorem~3.1]{longcompact}. A joint endpoint identity allows
an additional cubic factor while keeping all positive power moments
zero. Set
\[
 U=1+x,\qquad w=tU,\qquad
 e=\frac{U}{x}(1-w^2)=U\bigl(2y-(2+x)t^2\bigr),\qquad
 g=e+w^2=U(2y-t^2).
\]
For integers \(r\ge1\) and \(s\ge0\) put
\[
 J_{r,s}(V)=\frac12\int_{-V}^V z^s(1-z^2)^r\,dz.
\]
Here the endpoints are evaluated in a polynomial primitive.
Both endpoint integrands vanish to order at least \(r\) at \(V=1\),
so \(J_{r,s}(V)-J_{r,s}(1)\) is divisible by \((V-1)^{r+1}\).
The substitution \(z=tU\) gives
\[
 \Lambda(e^r w^s)=[x^r]U^{r-1}J_{r,s}(U)=0,\qquad
 \Lambda(xe^r w^s)=[x^{r-1}]U^{r-1}J_{r,s}(U)=J_{r,s}(1).
\]
In each coefficient the endpoint error is of order at least \(r+1\),
and the remaining factor \(U^{r-1}\) is monic of degree \(r-1\).
Thus in all polynomials divisible by \(e\) in \(K[e,w]\), the trace
vanishes, and the \(x\)-marked trace agrees with the interval
substitution \(e=1-z^2\), \(w=z\). In particular, for \(a,b\ge0\),
the phase \(e(1+w)^a g^b\) is pure. Its \(x\)-marked \(m\)th power
has trace
\begin{equation}\label{eq:spectrum-beta}
 \frac12\int_{-1}^1(1-z^2)^m(1+z)^{am}\,dz
 =2^{(a+2)m}\frac{((a+1)m)!\,m!}{((a+2)m+1)!}.
\end{equation}
The factor \(g^b\) disappears from this integral since \(g=e+w^2=1\)
under the marked substitution. Equation \eqref{eq:spectrum-beta}
follows by putting \(z=2u-1\) and integrating the beta monomial
\(u^{(a+1)m}(1-u)^m\).

The corresponding integer forms can be written using
\[
 L=R+Z,\quad D=2WR^2-(2R+Z)T^2,\quad E=2WR-T^2,\quad
 A_2=R^2+TL,\quad G_3=LE.
\]
For \(a,b\ge0\) set \(F_{a,b}=LD\,A_2^aG_3^b\) and \(d=4+2a+3b\).
Each factor is balanced homogeneous, and the sphere restriction is
\(e(1+w)^a g^b\). Equations \eqref{eq:spectrum-contraction} and
\eqref{eq:spectrum-beta} give
\[
 \mathcal E(F_{a,b}^m)=0,\qquad
 \mathcal E(ZF_{a,b}^m)=(dm+2)!\,2^{(a+2)m}
             \frac{((a+1)m)!\,m!}{((a+2)m+1)!}\ne0
 \qquad(m\ge1).
\]
At \((\xi_1,\xi_2,z_1,z_2)=(1,1,1,-1)\) we have \(R=0\) and
\(F_{a,b}=-4(-2)^a4^b\ne0\), which proves radial primitivity.
Specializing \(\xi_1=z_1=1\), \(z_2=0\), \(\xi_2=u\) yields
\(2^a(4u-2)(4u-1)^b\), with a simple zero at \(u=1/2\). Any scalar
multiple of a proper power would specialize to zero or to a scalar
multiple of a proper power. Hence these forms remain nonpowers over
every field extension. Taking \(a=b=0\) realizes degree four,
taking \(b=0\) realizes every even degree at least six, and taking
\(b=1\) realizes every odd degree at least seven.
